\section{Dataset and Vignette Design}
\label{sec:dataset}

We construct a factorially controlled set of relationship-dispute vignettes
designed to isolate the causal effect of a norm-violating partner's stated
gender on an LLM's fault-attribution judgment, holding every other narrative
element fixed.

\paragraph{Design.} The core confirmatory set crosses 9 relationship norm
families -- emotional labor, household labor, childcare, mental load,
financial provision, jealousy/possessiveness, sexuality and intimacy, career
sacrifice, and family obligations -- with 4 distinct scenarios per family (36
scenarios total), 4 gender configurations (male agent/female partner,
female agent/male partner, male-male, female-female), and 2 severity levels
(mild/severe), for $9 \times 4 \times 4 \times 2 = 288$ vignettes. Each
scenario is a distinct combination of task/object, violation form, and
obligation source within its family, giving within-family replication rather
than a single hand-written template per family. All 288 vignettes are
deterministically rendered from a single machine-readable design
specification, so that no cell is hand-written independently of the others.

\paragraph{Canonical structure.} Every vignette follows the same fixed
7-beat sentence order: (1) relationship context, (2) an explicitly stated
shared obligation, (3) the agent's norm-violating action at the assigned
severity level, (4) a concrete downstream consequence matched in specificity
across severity levels, (5) the agent's explanation, (6) the partner's
reaction, and (7) a fixed evaluation question. Holding this order and
structure fixed across all 288 vignettes ensures that models respond to the
manipulated conditions rather than to incidental presentation differences.

\paragraph{Obligation grounding.} Rather than leaving the source of the
agent's responsibility to be inferred, every vignette explicitly grounds it
in one of 8 standardized obligation-source types, each realized through a
fixed sentence template with only its content slots filled in, and each
tied to a specific literature basis rather than an ad hoc narrative choice.
\emph{Accepted role responsibility} (e.g.\ household labor, childcare) uses
Hardimon's distinction between role \emph{membership} and role
\emph{assignment} \citep{hardimon1994}: partnership or marriage alone does
not determine who does a given task, a specific accepted assignment does.
\emph{Established joint practice} (mental load) instead grounds the
obligation in an unspoken, reliance-generated pattern rather than a
promissory moment, following Daminger's account of mental load as often
unassigned and unspoken \citep{daminger2019}. \emph{Need-responsive
relational duty} (emotional labor) follows Clark and Mills' distinction
between communal, need-responsive relationships and transaction-tracking
exchange relationships \citep{clarkmills1979,clarkmills1993}.
\emph{Contribution-based reciprocity} (e.g.\ career sacrifice, family
obligations) grounds the obligation in fairness generated by unequal
contribution over time rather than a discrete promise, following
Rawls on fair play, Gouldner on the norm of reciprocity, and Walster,
Walster, and Berscheid's equity theory
\citep{rawls1964,gouldner1960,walster1978}. \emph{Recognized reliance on
disclosure} (sexuality and intimacy) draws on Scanlon's contractualist
account of promissory obligation, on which the wrong in frustrating
another's expectation lies in the assurance and induced reliance one has
knowingly created, not in the bare act of utterance \citep{scanlon1998}.
\emph{Baseline relational norm} (jealousy/possessiveness) grounds a duty
not to coerce or violate autonomy that holds independent of any relational
agreement \citep{scanlon1998,wertheimer1987}, and carries an explicit pilot
flag confirming its mild condition reads as an actual violation rather
than a reasonable boundary request. \emph{Good-faith relationship
maintenance} and \emph{fair notice of expectations} (both developed for
the sexuality and intimacy family) are general fairness principles without
a single named academic source, and are documented as such rather than
attributed to a citation they don't actually rest on. Obligation strength
is held constant across all cells within a family so that severity and
gender-configuration comparisons are not confounded by an incidentally
stronger or weaker stated obligation.

\paragraph{Fixed factors.} Two additional factors that could in principle
have been crossed are instead fixed within the core confirmatory design.
Intentionality is fixed at a single \emph{knowing-but-non-malicious} level
for every vignette -- the agent understood the obligation and chose not to
meet it, without acting out of malice -- rather than crossed across an
accidental/negligent/knowing/purposeful scale; a richer
intentionality-robustness arm using that fuller scale on 2 selected
scenarios per family is designed but deliberately deferred until the core
set has been piloted. Relationship context (relationship type and duration)
is likewise fixed per family at one canonical value (e.g.\ married with one
child for the childcare family; dating for the jealousy family) rather than
crossed with dating/married as a factor.

\paragraph{Agent labeling.} Agents are referred to by anonymized role labels
rather than names: \texttt{Agent 1} always denotes the norm-violating role
and \texttt{Agent 2} the partner role, fixed regardless of family or gender
configuration. Gender is conveyed exactly once per agent, as an explicit
tag at that agent's first mention (e.g.\ ``Agent 1 (female) and Agent 2
(male) have been cohabiting for 3 years...''); no pronouns appear anywhere
in the rendered text. This removes two confounds a named-agent design could
not fully rule out: residual non-gender connotation carried even by
gender-neutral names (cultural or ethnic association, familiarity,
frequency), and redundant double-encoding of gender through both name and
repeated pronoun use, which could differentially draw model attention to
the manipulated signal.

\paragraph{Evaluation framing.} Each vignette closes with the fixed question
``Was Agent 1 in the wrong?'' rather than the AITA subreddit's native ``the
asshole'' phrasing. This adopts the WHO-task reframing established by
Scruples \citep{lourie2021scruples}, the standard large-scale
AITA-derived benchmark for this class of task, rather than introducing a new,
untested wording -- avoiding the need for a separate wording-robustness
control arm to rule out interpretation differences.

\paragraph{Quality control.} All 288 vignettes were checked for word-count
parity within each scenario's 8 cells (no scenario exceeds a 15\% spread),
a banned-language scan (no moral adjectives, personality descriptors, or
gendered language outside the explicit tag), and a per-family
scenario-balance requirement (at least 2 distinct obligation sources, 4
distinct task/objects, and at least 3 distinct violation forms per family).
A tone-consistency audit verified that the fixed
knowing-but-non-malicious explanation reads comparably across all 9
families, correcting two families that initially under- or over-shot the
intended tone. A subsequent manual review of a 20-vignette stratified
sample surfaced a further class of issue specific to the mild/severe
design: fields shared verbatim across both severity cells (the
obligation sentence, the explanation, the partner response) can silently
imply an ongoing pattern that contradicts a mild cell's single-incident
framing, or leave a demonstrative reference ("this friendship") without
an antecedent. 17 of 36 scenarios were corrected for this, and the check
was subsequently automated (verb-aspect concordance, obligation-sentence
chronicity, and demonstrative-antecedent auditing) into a linter run
against every drafted scenario, converging to zero flagged issues.
